\section{Planarity}

% =========================================================
\begin{thmbox}{Theorem 1}
The graphs $K_5$ and $K_{3,3}$ are non-planar.
\end{thmbox}

\begin{proof}
Suppose first that $K_{3,3}$ is planar. Since $K_{3,3}$ has a cycle $u \to v \to w \to x \to y \to z \to u$ of length $6$, any planar drawing must contain this cycle drawn in the form of a hexagon.

Now the edge $wz$ must lie either wholly inside the hexagon or wholly outside it. We deal with the case in which $wz$ lies inside the hexagon --- the other case is similar. Since the edge $ux$ must not cross the edge $wz$, it must lie outside the hexagon. It is then impossible to draw the edge $vy$, as it would cross either $ux$ or $wz$.

\begin{center}
\begin{tikzpicture}[scale=1.0, every node/.style={circle,draw=accentblue,fill=accentblue!15,minimum size=6mm,inner sep=0pt}]
  \node (u) at (90:1.8)   {$u$};
  \node (v) at (30:1.8)   {$v$};
  \node (w) at (-30:1.8)  {$w$};
  \node (x) at (-90:1.8)  {$x$};
  \node (y) at (-150:1.8) {$y$};
  \node (z) at (150:1.8)  {$z$};

  \draw (u)--(v)--(w)--(x)--(y)--(z)--(u);
  \draw (w) -- (z);
  \draw (u) to[bend right=75] (x);
  \draw[highlightred, thick] (v) to[bend right=75] (y);
\end{tikzpicture}
\end{center}

Now suppose $K_5$ is planar. Since $K_5$ has a cycle $v \to w \to x \to y \to z \to v$ of length $5$, any planar drawing must contain this cycle drawn in the form of a pentagon.

Now the edge $wz$ must lie either wholly inside the pentagon or wholly outside it. We deal with the case in which $wz$ lies inside the pentagon --- the other case is similar. Since the edges $vx$ and $vy$ must not cross the edge $wz$, they must both lie outside the pentagon. But the edge $xz$ cannot cross $vy$, and so must lie inside the pentagon; similarly, $wy$ cannot cross $vx$, and so must also lie inside. Then $wy$ and $xz$ must cross, a contradiction.

\begin{center}
\begin{tikzpicture}[scale=1.0, every node/.style={circle,draw=accentblue,fill=accentblue!15,minimum size=6mm,inner sep=0pt}]
  \node (v) at (90:1.8)   {$v$};
  \node (w) at (18:1.8)   {$w$};
  \node (x) at (-54:1.8)  {$x$};
  \node (y) at (-126:1.8) {$y$};
  \node (z) at (162:1.8)  {$z$};

  \draw (v)--(w)--(x)--(y)--(z)--(v);
  \draw (w) -- (z);
  \draw (x) -- (z);
  \draw[highlightred, thick] (w) -- (y);
  \draw (v) to[bend right=35] (x);
  \draw (v) to[bend left=35] (y);
\end{tikzpicture}
\end{center}
\end{proof}

% =========================================================
\begin{thmbox}{{Theorem 2 (Kuratowski, 1930)}}
A graph is planar if and only if no subgraph is homeomorphic to $K_5$ or $K_{3,3}$.
\end{thmbox}

% =========================================================
\begin{thmbox}{Theorem 3}
A graph is planar if and only if no subgraph is contractible to $K_5$ or $K_{3,3}$.
\end{thmbox}

% =========================================================
\begin{thmbox}{Handshaking Lemma (for Faces)}
In any plane drawing of a planar graph,
$$\sum \deg(f_i) = 2\,|E(G)|.$$
\end{thmbox}

\begin{proof}
In a plane drawing, each edge has two sides (two incident faces):
\begin{itemize}
  \item the edge may lie on the boundary of a single face (counted twice), or
  \item the edge lies on the boundary of a face $f_1$ on one side, and a face $f_2$ on the other side.
\end{itemize}
In either case, each edge contributes exactly $2$ to the total degree. Hence $\sum \deg(f_i) = 2\,|E(G)|$.
\end{proof}

% =========================================================
\begin{thmbox}{{Theorem 4 (Euler, 1750)}}
Let $G$ be a plane drawing of a connected planar graph, and let $n$, $m$, and $F$ denote respectively the number of vertices, edges, and faces of $G$. Then
$$n - m + F = 2.$$
\end{thmbox}

\begin{proof}
By induction on the number of edges of $G$. If $m = 0$, then $n = 1$, $F = 1$. The theorem is true in this case.

Now, suppose that the theorem holds for all graphs with at most $m-1$ edges. Let $G$ be a graph with $m$ edges. If $G$ is a tree, then $m = n-1$, $F = 1$, so that indeed $n - m + F = 2$ as required.

If $G$ is not a tree, let $e$ be an edge in some cycle of $G$. Then $G - e$ is a connected plane graph with $n$ vertices, $m-1$ edges, $F-1$ faces, so that $n - (m-1) + (F-1) = 2$ by the induction hypothesis. Adding back the removed edge gives $n - m + F = 2$, as required.
\end{proof}

% =========================================================
\begin{thmbox}{Corollary 1}
\begin{enumerate}
  \item[(i)] If $G$ is a connected simple planar graph with $n\ (\geq 3)$ vertices and $m$ edges, then $m \leq 3n - 6$.
  \item[(ii)] If, in addition, $G$ has no triangle, then $m \leq 2n - 4$.
\end{enumerate}
\end{thmbox}

\begin{proof}
(i): Since each face is bounded by at least three edges, it follows, on counting up the edges around each face, that
$$3F \leq 2m$$
(the factor $2$ appears since each edge bounds two faces). We obtain the required bound by replacing $F = 2 - n + m$.

\medskip
\noindent (ii): This follows in a similar way, by replacing $3F \leq 2m$ with $4F \leq 2m$.
\end{proof}

\medskip
\noindent As a consequence of Corollary 1, we obtain a second, shorter proof that $K_5$ and $K_{3,3}$ are non-planar (Theorem 1).

\begin{proof}[Proof 2 of Theorem 1 (via Euler's Formula)]
$K_{3,3}$: $n=6$, $m=9$ (no triangles). Since $2n-4 = 8 < m = 9$, $K_{3,3}$ is non-planar.

\medskip
\noindent $K_5$: $n=5$, $m=10$. Since $3n-6 = 9 < m = 10$, $K_5$ is non-planar.
\end{proof}

% =========================================================
\begin{thmbox}{Corollary 2}
Let $G$ be a simple connected planar graph. Then $G$ contains a vertex of degree $5$ or less.
\end{thmbox}

\begin{proof}
Suppose, for contradiction, that $d_i \geq 6$ for every vertex. Then
$$\sum d_i \geq 6n \implies 2m \geq 6n \implies m \geq 3n > 3n - 6,$$
contradicting $m \leq 3n-6$ (Corollary 1(i)).
\end{proof}
